\section{Introduction}

Insider threat detection within high-security and tiered-clearance organizational networks remains one of the most intractable challenges in modern cybersecurity. Unlike external cyber-attackers who must breach peripheral firewalls and intrusion detection systems, malicious insiders possess legitimate credentials, authorized system privileges, and familiar access patterns. Recent security intelligence reports indicate that espionage-grade insider incidents—characterized by low-volume, slow-burn data collection and deliberate metadata manipulation—incur an average organizational remediation cost exceeding \$16.2 million per breach \citep{wei2024ewatcher, xiao2024sentinel}. Conventional security information and event management (SIEM) solutions and single-layer machine learning classifiers predominantly rely on point-in-time endpoint logs or crude volume-threshold alerts. However, out-of-distribution (OOD) adversaries routinely exploit these systems through mimicry tactics, operating during standard business hours, blending with benign traffic, or stripping file header metadata prior to exfiltration \citep{shoderu2025dfrbust, ga02023dtgi}. Under such adversarial conditions, single-layer statistical detectors suffer severe performance degradation, with ROC-AUC dropping by up to 26.25\% when confronted with credential-phishing and author-spoofing evasion vectors.

Beyond the primary challenge of real-time detection, security operations centers (SOCs) face a second critical vulnerability: the forensic evidence delivery gap. Standard behavioral anomaly alerts generated by machine learning models are typically uncalibrated black-box outputs that fail to satisfy judicial standards of admissibility. When a breach is escalated to legal prosecution or counter-intelligence enforcement, digital evidence must demonstrate strict compliance with the ISO/IEC 27043 international standard for digital forensic readiness \citep{mavroeidis2018ps0}. Existing threat detection architectures generate ephemeral log trails that lack cryptographic tamper-evidence, chain-of-custody verification, and formal semantic provenance alignment (such as W3C-PROV). Consequently, forensic investigators are forced to manually perform costly, retrospective log aggregation that delays response times and risks evidence contamination.

To bridge the gap between unstructured endpoint telemetry and court-admissible forensic evidence, an under-explored signal resides within cross-format file metadata and spatio-temporal provenance. Enterprise document formats—including Office Open XML (DOCX), Portable Document Format (PDF), Exchangeable Image File Format (EXIF), Electronic Mail (EML), and Git version-control repositories—embed rich structural provenance metadata (e.g., revision identifiers, author history, editing timestamps, and equipment signatures). When systematically extracted and mapped into a graph representation, these cross-format metadata attributes form a dense provenance knowledge graph. This temporal graph structure exposes behavioral anomalies and unauthorized clearance probing that are completely invisible to isolated log parsers \citep{umair2024exif2vec, yasenenko2025imagemeta}.

In this paper, we introduce \textbf{TURRET OS}, an integrated, multi-layered counter-espionage framework that unifies cross-format document metadata harvest, a W3C-PROV-aligned provenance knowledge graph, a boolean rule Abstract Syntax Tree (AST) engine, a spatio-temporal Graph Neural Network (GNN), and Ed25519-signed Merkle evidence packaging. TURRET OS addresses four fundamental research questions:
\begin{itemize}
    \item \textbf{RQ1 (Detection Accuracy)}: Can a multi-layered provenance graph architecture achieve state-of-the-art detection performance (ROC-AUC $\ge 0.99$, F1 $\ge 0.88$) on low-prevalence espionage datasets while maintaining zero-delay operational response ($<0.1$ min)?
    \item \textbf{RQ2 (Ablation \& Leakage)}: Does graph structural topology provide authentic discriminative signal rather than feature leakage across temporal splits?
    \item \textbf{RQ3 (Adversarial Robustness)}: Can joint multi-layer inference mitigate OOD evasion tactics (mimicry, metadata-stripping, identity-proxy spoofing) that cause single-layer detectors to fail?
    \item \textbf{RQ4 (Forensic Readiness)}: Can automated detection pipelines generate cryptographically verifiable evidence bundles satisfying 100\% of ISO/IEC 27043 readiness attributes?
\end{itemize}

The key contributions of TURRET OS are five-fold:
\begin{enumerate}
    \item \textbf{Multi-Layer Joint Architecture}: We present a 4-layer braided pipeline integrating L1 multi-parser metadata harvest, L2 Neo4j provenance graph construction, L3 boolean rule AST and GNN inference, and L4 cryptographically signed evidence packaging.
    \item \textbf{High Detection Performance}: On the 141,725-record D3 TSEC synthetic corpus, TURRET OS achieves an uncalibrated ROC-AUC of 0.9983, a calibrated best F1-score of 0.9747 (at threshold 0.37), an F1-score of 0.8813 at a strict 0.5\% FPR operating point, an MCC of 0.9732, and an operational time-to-detect of $<0.1$ minutes.
    \item \textbf{Adversarial Red-Team Resilience}: We demonstrate that while single-layer L1 feature detectors degrade to 0.7363 ROC-AUC (-26.25\% drop) under OOD credential-phishing attacks, TURRET OS's multi-layered graph topology and rule AST compensate for feature evasion.
    \item \textbf{100\% ISO/IEC 27043 Forensic Readiness}: TURRET OS automatically packages detection alerts into Ed25519-signed ZIP evidence bundles containing Merkle-tree root hashes, achieving 100.0\% coverage across all eight ISO/IEC 27043 forensic attributes.
    \item \textbf{Reproducibility \& Runtime Safety}: We provide a fully reproducible codebase, wall-clock hardware profile (80 CPU cores, 503.58 GB RAM, NVIDIA L40S), and a deployed SHAP feature-whitelist runtime safety invariant that enforces zero label leakage during enterprise deployment.
\end{enumerate}
